\documentclass{article}

\usepackage[margin=1in]{geometry}
\usepackage{amsmath,amsthm,amssymb}
\usepackage[spanish]{babel}
\usepackage{enumitem}

\theoremstyle{plain}
\newtheorem{proposicion}{Proposición}[section]

\theoremstyle{definition}
\newtheorem{definition}{Definición}[section]

% Number sets
\newcommand{\R}{\mathbb{R}}
\newcommand{\Z}{\mathbb{Z}}
\newcommand{\N}{\mathbb{N}}
\newcommand{\Q}{\mathbb{Q}}
\newcommand{\C}{\mathbb{C}}

% Topology operators
\newcommand{\cl}[1]{\overline{#1}}              % Closure
\newcommand{\Int}[1]{{#1}^\circ}                % Interior
\newcommand{\bd}{\partial}                      % Boundary
\newcommand{\Ent}[1]{\mathcal{O}(#1)}           % Neighborhoods (Entornos)
\newcommand{\pf}{\mathcal{P}_F}                 % Finite parts
\newcommand{\partes}{\mathcal{P}}               % Power set

% Useful shortcuts
\newcommand{\sii}{\iff}                          % Si y sólo si
\newcommand{\imp}{\implies}                      % Implies
\newcommand{\setdiff}{\smallsetminus}            % Set difference

% Net convergence/accumulation notation
\newcommand{\evto}[1]{\stackrel{E}{\hookrightarrow} #1}   % Eventually in
\newcommand{\frto}[1]{\stackrel{F}{\hookrightarrow} #1}   % Frequently in
\newcommand{\nnevto}[1]{\not\!\stackrel{E}{\hookrightarrow} #1}
\newcommand{\nnfrto}[1]{\not\!\stackrel{F}{\hookrightarrow} #1}

% Exercise environment
\newenvironment{ejercicio}[1]{\begin{trivlist}
\item[\hskip \labelsep {\bfseries Ejercicio}\hskip \labelsep {\bfseries #1.}]}{\end{trivlist}}

\setlist[enumerate,1]{label=(\alph*), leftmargin=*, itemsep=2pt}

\begin{document}

\title{Topología: Práctica V}
\author{Bustos Jordi\\Compactos}

\maketitle

\section*{Compactos}

\begin{ejercicio}{1}
    Demostrar que si X es un espacio topológico compacto y Hausdorff bajo dos topologías y entonces \( \tau=\sigma \) o bien no son comparables entre sí.
\end{ejercicio}
\begin{proof}
    Supongamos que las topologías \( \tau \) y \( \sigma \) son comparables. Sin pérdida de generalidad, podemos asumir que una está contenida en la otra; supongamos que \( \sigma \subseteq \tau \).

    Consideremos la aplicación identidad:
    \[ id: (X, \tau) \longrightarrow (X, \sigma) \]
    definida por \( id(x) = x \) para todo \( x \in X \).

    A partir de las hipótesis de nuestro problema, observamos lo siguiente:
    \begin{enumerate}
        \item \( id \) es claramente una función inyectiva (y sobreyectiva).
        \item \( id \) es continua, ya que para cualquier abierto \( U \in \sigma \), su preimagen es \( id^{-1}(U) = U \), y como por suposición \( \sigma \subseteq \tau \), se cumple que \( U \in \tau \).
        \item El espacio de partida \( (X, \tau) \) es compacto por hipótesis.
        \item El espacio de llegada \( (X, \sigma) \) es Hausdorff por hipótesis.
    \end{enumerate}

    Utilizando la \textbf{Proposición 6.2.7} del apunte de Demetrio (si \( f: K \to X \) es continua e inyectiva con \( K \) compacto y \( X \) Hausdorff, entonces \( f: K \to f(K) \) es un homeomorfismo), y dado que la imagen de nuestra función es todo el espacio de llegada (\( id(X) = X \)), concluimos que:
    \[ id: (X, \tau) \longrightarrow (X, \sigma) \]
    es un homeomorfismo.

    Al ser un homeomorfismo, \( id \) es una aplicación abierta. Esto significa que si tomamos cualquier conjunto abierto \( V \in \tau \), su imagen \( id(V) = V \) debe ser un abierto en \( \sigma \), lo que demuestra la inclusión inversa \( \tau \subseteq \sigma \).

    Como ya teníamos que \( \sigma \subseteq \tau \) y ahora hemos probado que \( \tau \subseteq \sigma \), se deduce que \( \tau = \sigma \). Por lo tanto, si las topologías son comparables entre sí, obligatoriamente deben ser iguales, lo que demuestra que \( \tau = \sigma \) o bien no son comparables.
\end{proof}

\clearpage

\begin{ejercicio}{2 (Lema del tubo)}
    Sean X, Y espacios topológicos y supongamos que Y es compacto. Sea U un abierto de \( X \times Y \) (respecto de la topología producto) que contiene a la recta vertical \( x_{0}\times Y \). Demostrar que \( \exists W \in \mathcal{O}_{X}^{a}(x_{0}) \) tal que \( W \times Y \subset U \).
\end{ejercicio}
\begin{proof}
    Dado \( U \in \tau_P \) con \( \{ x_0 \} \times Y \subseteq U \), tenemos que para cada punto \( (x_0, y) \in \{ x_0 \} \times Y \) existen dos abiertos \( U_{y} \in \mathcal{O}_X^a(x_0) \) y \( V_{y} \in \mathcal{O}_Y^a(y) \) tales que \( (x_0, y) \in U_{y} \times V_{y} \subseteq U \).

    La familia de conjuntos \( {\{ V_{y} \}}_{y \in Y} \) forma un cubrimiento abierto para \( Y \). Como \( Y \) es compacto, podemos extraer un subcubrimiento finito, es decir, existen \( y_1, \ldots, y_n \in Y \) tales que \( Y \subseteq \bigcup_{i = 1}^n V_{y_i} \).

    Definimos \( W = \bigcap_{i = 1}^n U_{y_i} \). Como \( W \) es una intersección finita de abiertos en \( X \) que contienen a \( x_0 \), resulta que \( W \in \mathcal{O}_X^a(x_0) \).

    Finalmente, veamos que cumple \( W \times Y \subseteq U \). Sea \( (x, y) \in W \times Y \). Como los \( V_{y_i} \) cubren a \( Y \), existe algún \( i \in \{1, \ldots, n\} \) tal que \( y \in V_{y_i} \). Por otro lado, como \( x \in W \) y \( W \) es la intersección de todos los \( U_{y_j} \), en particular se tiene que \( x \in U_{y_i} \).

    Por lo tanto, \( (x, y) \in U_{y_i} \times V_{y_i} \subseteq U \). Como esto vale para cualquier punto de \( W \times Y \), concluimos que \( W \times Y \subset U \).
\end{proof}

\begin{ejercicio}{3}
    Sea \( \mathcal{F} \) una familia de conjuntos con la PIF en \( X \).
    \begin{enumerate}
        \item[(a)] Dada \( f : X \rightarrow Y \) mostrar que la familia \( f(\mathcal{F}) \) tiene la PIF en Y.
        \item[(b)] Existe una familia \( M \) en \( X \) con la PIF, tal que \( \mathcal{F}\subseteq\mathcal{M} \) y \( \mathcal{M} \) maximal.
        \item[(c)] Si \( \mathcal{F} \) es maximal, entonces
              \begin{enumerate}
                  \item[i.] Si \( A \) y \( B \) pertenecen a \( \mathcal{F} \) entonces \(A \cap B \) pertenece a \( \mathcal{F} \).
                  \item[ii.] Si \( A \cap F \ne\varnothing \) para todo \(F\in\mathcal{F}\) entonces \(A\in\mathcal{F}\).
              \end{enumerate}
    \end{enumerate}
\end{ejercicio}
\begin{proof}
    Sea \( \mathcal{F} = {\{ F_i \}}_{i \in I} \) con la PIF en \( X \).
    \begin{enumerate}
        \item Tenemos que dado \( \mathbb{F} \subseteq I \) finito, vale que \( \bigcap_{i \in \mathbb{F}} X \cap F_i \neq \varnothing \) por lo que \( \exists a \in X \cap F_i \quad \forall i \in \mathbb{F} \).
              Esto nos dice que, dado \( \bigcap_{i \in \mathbb{F}} Y \cap f(F_i) \neq \varnothing \) tomando \( f(a) \therefore {\{ f(F_i) \}}_{i \in I} \) tiene la PIF en \( Y \).
        \item Definamos el siguiente conjunto de familias:
              \[ P := \{ \mathcal{A} \subseteq \mathcal{P}(X) : \mathcal{F} \subseteq \mathcal{A} \text{ y } \mathcal{A} \text{ tiene la PIF} \} \]

              Ordenado por la inclusión (\( \subseteq \)). Primero, notemos que \( \mathcal{F} \in P \neq \varnothing \).

              Para aplicar el Lema de Zorn, debemos verificar que toda cadena en \( P \) tiene una cota superior en \( P \).
              Sea \( C = {\{ \mathcal{A}_i \}}_{i \in I} \) una cadena totalmente ordenada en \( P \). Proponemos como cota superior a la unión de todos los elementos de la cadena:
              \[ \mathcal{U} = \bigcup_{i \in I} \mathcal{A}_i \]

              Veamos que \( \mathcal{U} \in P \):
              \begin{enumerate}
                  \item Es claro que \( \mathcal{F} \subseteq \mathcal{U} \).
                  \item Veamos que \( \mathcal{U} \) tiene la PIF en \( X \). Sean \( A_1, A_2, \ldots, A_n \in \mathcal{U} \).
                        Por definición de la unión, cada \( A_k \) pertenece a alguna familia \( \mathcal{A}_{i_k} \) de la cadena \( C \).
                        Como \( C \) es totalmente ordenada, dados finitos elementos de la cadena, siempre existe uno que los contiene a todos (el máximo entre ellos respecto a la inclusión).
                        Sea \( \mathcal{A}^* \in C \) esa familia tal que \( A_1, \ldots, A_n \in \mathcal{A}^* \).

                        Como \( \mathcal{A}^* \in P \), sabemos que \( \mathcal{A}^* \) tiene la PIF en \( X \). Por lo tanto, la intersección de estos finitos conjuntos es no vacía:
                        \[ \bigcap_{k=1}^n A_k \neq \varnothing \]
                        Esto prueba que \( \mathcal{U} \) tiene la PIF en \( X \).
              \end{enumerate}

              Al tener \( \mathcal{U} \in P \) y ser \( \mathcal{A}_i \subseteq \mathcal{U} \) para todo \( i \in I \), resulta que \( \mathcal{U} \) es cota superior de la cadena \( C \).

              Como toda cadena en \( P \) está acotada superiormente, por el Lema de Zorn, el conjunto \( P \) posee un elemento maximal \( \mathcal{M} \).

              Por pertenecer a \( P \), \( \mathcal{M} \) es una familia con la PIF que contiene a \( \mathcal{F} \). Por ser maximal en \( P \), es una familia maximal con la PIF (si existiera una familia más grande con la PIF, esta también contendría a \( \mathcal{F} \), estaría en \( P \) y contradiría la maximalidad de \( \mathcal{M} \)).
        \item Supongamos ahora que \( \mathcal{F} \) es maximal. Es fácil ver que si \( A, B \in \mathcal{F} \) y \( A \cap B \notin \mathcal{F} \) entonces \( \mathcal{F} \cup \{ A \cap B \} \) tiene la PIF y, por lo tanto, \( \mathcal{F} \) no es maximal.
              Para la otra procedemos exactamente igual, si \( A \notin \mathcal{F} \) vale que \( \mathcal{F} \subseteq \{ A \} \cup \mathcal{F} \) y es fácil ver que esta extensión tiene la PIF en \( X \), por lo tanto \( \mathcal{F} \) no es maximal.
    \end{enumerate}
\end{proof}

\begin{ejercicio}{4}
    Dados X e Y espacios topológicos con Y compacto. Entonces la proyección \( \pi_{x}:X\times Y\rightarrow X \) es cerrada.
\end{ejercicio}
\begin{proof}
    Sea \( C \subseteq X \times Y \) cerrado y una red \( {\{x_i\}}_{i \in I} \subseteq \pi_x(C) \) tal que \( x_i \longrightarrow x \) en \( X \). Queremos ver que \( x \in \pi_x(C) \).

    Dado que cada \( x_i \in \pi_x(C) \), para cada \( i \in I \) existe al menos un elemento \( y_i \in Y \) tal que \( (x_i, y_i) \in C \).
    Esto nos permite construir una red \( {\{ y_i \}}_{i \in I} \) en \( Y \).

    Como \( Y \) es compacto por hipótesis, existe una subred \( {\{y_{i_j}\}}_{j \in J} \) de \( {\{y_i\}}_{i \in I} \) tal que \( y_{i_j} \longrightarrow y \in Y \).

    Ahora consideramos la subred correspondiente en el espacio producto \( X \times Y \), dada por \( {\{(x_{i_j}, y_{i_j})\}}_{j \in J} \).

    Como la red original \( {\{x_i\}}_{i \in I} \) converge a \( x \), cualquier subred suya también debe converger al mismo límite, lo que nos asegura que \( x_{i_j} \longrightarrow x \).

    Por las propiedades de la topología producto, una red de pares converge si y solo si converge en cada una de sus coordenadas por separado.
    Como \( x_{i_j} \longrightarrow x \) e \( y_{i_j} \longrightarrow y \), se deduce que:
    \[ (x_{i_j}, y_{i_j}) \longrightarrow (x, y) \quad \text{en } X \times Y \]

    Puesto que \( (x_{i_j}, y_{i_j}) \in C \) para todo \( j \in J \) y el conjunto \( C \) es cerrado, el límite de esta subred también debe pertenecer a \( C \).
    De este modo, concluimos que \( (x, y) \in C \).

    Finalmente, aplicando la proyección a este punto límite, obtenemos:
    \[ \pi_x(x, y) = x \in \pi_x(C) \]

    Esto demuestra que \( \pi_x(C) \) contiene a todos sus puntos de acumulación y, por lo tanto, es un conjunto cerrado en \( X \).
\end{proof}

\begin{ejercicio}{5}
    Sea \( f : X \to Y \) donde \( X \) e \( Y \) son espacios topológicos con \( Y \) Hausdorff. Ya hemos visto (práctica 2) que \( f \in C(X,Y) \implies Gr_f = \{ (x, f(x)) : x \in X\} \) es cerrado en \( X \times Y \).

    Probar que la recíproca es cierta si Y es compacto.
\end{ejercicio}
\begin{proof}
    Supongamos que \( Y \) es un espacio compacto y Hausdorff, y que el gráfico de la función, \( Gr_f \), es cerrado en \( X \times Y \).
    Queremos probar que \( f \) es continua.

    Para ello, veamos que para cada cerrado \( F \subseteq Y \), \( f^{-1}(F) \) es cerrada en \( X \).

    Sea \( F \subseteq Y \) un conjunto cerrado cualquiera. Consideremos las proyecciones naturales del espacio producto:
    \[ \pi_x : X \times Y \longrightarrow X \quad \text{y} \quad \pi_y : X \times Y \longrightarrow Y \]

    Observamos que el conjunto \( X \times F \) es cerrado en \( X \times Y \) por definición de la topología producto.

    Como la intersección de dos conjuntos cerrados es cerrada, el conjunto:
    \[ Gr_f \cap (X \times F) \]
    es un cerrado en \( X \times Y \). Notemos qué elementos componen esta intersección:
    \[ Gr_f \cap (X \times F) = \{ (x, f(x)) \in X \times Y : f(x) \in F \} \]

    Ahora aplicamos la proyección en la primera coordenada, \( \pi_x \), a dicho conjunto.
    Por el \textbf{Ejercicio 4}, dado que \( Y \) es compacto, sabemos que la proyección \( \pi_x \) es una aplicación cerrada.
    Por lo tanto, la imagen de nuestro cerrado bajo \( \pi_x \) debe ser un cerrado en \( X \):
    \[ \pi_x \big( Gr_f \cap (X \times F) \big) = \{ x \in X : f(x) \in F \} = f^{-1}(F) \]

    Como \( f^{-1}(F) \) es la imagen de un cerrado bajo una aplicación cerrada, concluimos que \( f^{-1}(F) \) es un conjunto cerrado en \( X \).

    Al ser \( F \) un cerrado arbitrario de \( Y \), queda demostrado que \( f \) es una función continua.
\end{proof}

\begin{ejercicio}{6}
    Dado un \( X \) cualquiera dotado de la topología co-finita, entonces \( (X, \tau_{CF}) \) es compacto.
\end{ejercicio}
\begin{proof}
    Recordemos que la topología co-finita en \( X \) se define como:
    \[ \tau_{CF} = \{ U \subseteq X \, : \, X \setminus U \text{ es finito} \} \cup \{ \varnothing \} \]

    Sea \( {\{ U_i \}}_{i \in I} \) un cubrimiento abierto arbitrario de \( X \). Si \( X = \varnothing \), el espacio es trivialmente compacto. Supongamos entonces que \( X \neq \varnothing \).

    Como los \( U_i \) cubren a \( X \), debe existir al menos un conjunto en el cubrimiento que sea no vacío. Elijamos uno de ellos y llamémoslo \( U_{i_0} \), con \( i_0 \in I \).

    Por pertenecer a \( \tau_{CF} \) y ser distinto de vacío, por definición sabemos que su complemento \( X \setminus U_{i_0} \) es un conjunto finito. Supongamos entonces que tiene \( n \) elementos:
    \[ X \setminus U_{i_0} = \{ x_1, x_2, \dots, x_n \} \]

    Dado que la familia original \( {\{ U_i \}}_{i \in I} \) cubría a todo \( X \), en particular cubre a estos \( n \) puntos. Por lo tanto, para cada \( k \in \{1, \dots, n\} \), podemos elegir un índice \( i_k \in I \) tal que \( x_k \in U_{i_k} \).

    Consideremos ahora la subfamilia formada por \( U_{i_0} \) y los conjuntos que acabamos de elegir:
    \[ \{ U_{i_0}, U_{i_1}, U_{i_2}, \dots, U_{i_n} \} \]

    Esta familia finita es un subcubrimiento de \( X \). En efecto, si tomamos cualquier \( x \in X \), tenemos dos opciones:
    \begin{itemize}
        \item Si \( x \in U_{i_0} \), ya está cubierto por \( U_{i_0} \).
        \item Si \( x \notin U_{i_0} \), entonces \( x \in X \setminus U_{i_0} \), por lo que \( x = x_k \) para algún \( k \), y en consecuencia \( x \in U_{i_k} \).
    \end{itemize}

    Por lo tanto:
    \[ X \subseteq U_{i_0} \cup \left( \bigcup_{k=1}^n U_{i_k} \right) \]

    Como de un cubrimiento abierto cualquiera hemos podido extraer un subcubrimiento finito, concluimos que \( (X, \tau_{CF}) \) es compacto.
\end{proof}

\begin{ejercicio}{7}
    Demostrar que \( A \) es totalmente acotado \( \iff \cl{A} \) es totalmente acotado.
\end{ejercicio}
\begin{proof}
    La vuelta es sencilla, si \( \forall \varepsilon > 0 \) vale que \( \exists x_1, \ldots, x_n \in X \, : \, \cl{A} \subseteq \bigcup_{i = 1}^n B_\varepsilon(x_i) \), como \( A \subseteq \cl{A} \) tomamos el mismo cubrimiento.

    Para la ida supongamos que \( A \) es un subconjunto totalmente acotado. Sea \( \varepsilon > 0 \) queremos probar que \( \cl{A} \) puede ser cubierto por una cantidad finita de bolas de radio \( \varepsilon \).

    Como \( A \) es totalmente acotado, para el radio \( \frac{\varepsilon}{2} \) existen \( x_1, \dots, x_n \in X \) tales que:
    \[ A \subseteq \bigcup_{i=1}^n B_{\frac{\varepsilon}{2}}(x_i) \]

    Sea \( y \in \cl{A} \) un punto cualquiera. Por definición de clausura, todo entorno de \( y \) interseca a \( A \).
    En particular, la bola \( B_{\frac{\varepsilon}{2}}(y) \) interseca a \( A \), lo que significa que existe un punto \( z \in A \) tal que:
    \[ d(y, z) < \frac{\varepsilon}{2} \]

    Como \( z \in A \), sabemos que \( z \) debe pertenecer a alguna de las bolas de nuestro subcubrimiento finito original. Por lo tanto, existe un índice \( i \in \{1, \dots, n\} \) tal que:
    \[ z \in B_{\frac{\varepsilon}{2}}(x_i) \implies d(z, x_i) < \frac{\varepsilon}{2} \]

    Utilizando la desigualdad triangular, podemos acotar la distancia entre \( y \) y el centro \( x_i \):
    \[ d(y, x_i) \leq d(y, z) + d(z, x_i) < \frac{\varepsilon}{2} + \frac{\varepsilon}{2} = \varepsilon \]

    Esto nos dice que \( y \in B_{\varepsilon}(x_i) \). Como \( y \) era un punto arbitrario de \( \cl{A} \), hemos demostrado que todo punto de la clausura está a distancia menor a \( \varepsilon \) de alguno de los \( n \) centros originales. Es decir:
    \[ \cl{A} \subseteq \bigcup_{i=1}^n B_{\varepsilon}(x_i) \]

    Como para todo \( \varepsilon > 0 \) hemos encontrado un cubrimiento finito por bolas de radio \( \varepsilon \), concluimos que \( \cl{A} \) es totalmente acotado.
\end{proof}

\begin{definition}
    Sea \( (X, d) \) un espacio métrico. Un número \( \delta > 0 \) se llama \textbf{número de Lebesgue} para un cubrimiento abierto \( \mathcal{U} \) de \( X \) si para cada subconjunto \( A \subseteq X \) con diámetro menor a \( \delta \), existe un conjunto \( U \in \mathcal{U} \) tal que \( A \subseteq U \).
\end{definition}

\begin{ejercicio}{8}
    Dar un ejemplo de 2 abiertos (con la topología usual) que cubran \( (0, 1) \) y que no tengan un \( \delta \) de Lebesgue.
\end{ejercicio}
\begin{proof}
    Definimos la sucesión de puntos \( x_n = \frac{1}{n} \) para \( n \ge 2 \). Notemos que el conjunto \( C = \left \{ \frac{1}{n} : n \ge 2 \right \} \) es cerrado en la topología de subespacio de \( (0, 1) \), ya que su único punto de acumulación en \( \mathbb{R} \) es el \( 0 \), que no pertenece a \( (0, 1) \).

    Definimos nuestro primer abierto \( U \) como todo el espacio salvo esos puntos:
    \[ U = (0, 1) \setminus C \]

    Para cubrir los puntos de \( C \), definimos nuestro segundo abierto \( V \) como una unión infinita de intervalitos centrados en cada \( \frac{1}{n} \) (nuestros solapamientos), pero con radios que se achican muy rápidamente, por ejemplo, al cubo:
    \[ V = \bigcup_{n=2}^\infty \left( \frac{1}{n} - \frac{1}{n^3}, \frac{1}{n} + \frac{1}{n^3} \right) \cap (0, 1) \]

    Como \( V \) es unión de intervalos abiertos, es abierto. Además, es claro que \( U \cup V = (0, 1) \), puesto que \( U \) cubre todo \( (0, 1) \) excepto los puntos de \( C \), y \( V \) cubre explícitamente cada punto de \( C \). Por lo tanto, \( \{U, V\} \) es un cubrimiento de 2 abiertos de \( (0, 1) \).

    Veamos por absurdo que no tiene número de Lebesgue. Supongamos que existe un \( \delta > 0 \) tal que todo conjunto de diámetro menor a \( \delta \) está completamente contenido en \( U \) o completamente contenido en \( V \).

    Por la propiedad arquimediana, podemos elegir un \( N \in \mathbb{N} \) lo suficientemente grande tal que:
    \[ \frac{2}{N^3} < \delta \]
    (Además, tomamos \( N \) lo bastante grande para que los intervalos de \( V \) no se intersequen entre sí, lo cual ocurre para todo \( N \ge 3 \)).

    Consideremos el conjunto \( A \) dado por un pequeño intervalo cerrado que contiene un punto de \( C \) y se excede apenas del solapamiento:\[
        A = \left[ \frac{1}{N} - \frac{2}{N^3}, \frac{1}{N} \right]
    \]

    Calculamos el diámetro de \( A \):
    \[ \text{diam}(A) = \frac{2}{N^3} < \delta \]

    Según la definición del número de Lebesgue, \( A \) debería estar contenido enteramente en \( U \) o en \( V \). Analicemos ambas opciones:
    \begin{enumerate}
        \item \( A \subseteq U \): No, porque el punto \( \frac{1}{N} \in A \), pero por construcción de nuestro primer abierto, \( \frac{1}{N} \notin U \).
        \item \( A \subseteq V \): No, el único componente de \( V \) que contiene al punto \( \frac{1}{N} \) es el intervalo \( \left( \frac{1}{N} - \frac{1}{N^3}, \frac{1}{N} + \frac{1}{N^3} \right) \). Sin embargo, el extremo inferior de \( A \) es \( \frac{1}{N} - \frac{2}{N^3} \), que es estrictamente menor al límite del intervalo \( \frac{1}{N} - \frac{1}{N^3} \). Por ende, \( A \) no cabe dentro de \( V \).
    \end{enumerate}

    Hemos encontrado un conjunto \( A \) con diámetro menor a \( \delta \) que no está contenido en ninguno de los dos abiertos (siempre necesitamos los dos para cubrirlo). Esto contradice que \( \delta \) sea un número de Lebesgue.

    Por lo tanto, el cubrimiento \( \{U, V\} \) no posee número de Lebesgue.
\end{proof}

\begin{definition}[Liptschitz continua]
    Una función \( f : X \to Y \) es Liptschitz continua \( \iff \) existe \( L > 0 \) tal que para todo \( x_1, x_2 \in X \) vale que \( d_Y(f(x_1), f(x_2)) \leq L \cdot d_X(x_1, x_2) \).
\end{definition}
\begin{definition}[Uniformemente continua]
    Una función \( f : X \to Y \) es uniformemente continua \( \iff \) \( \forall \varepsilon > 0 \) existe \( \delta > 0 \) tal que \( \forall x_1, y_1 \in X \) se cumple que si \( d_X(x_1, x_2) < \delta \implies d_Y(f(x_1), f(x_2)) < \varepsilon \).
\end{definition}

\begin{ejercicio}{9}
    Dados \( X \) e \( Y \) espacios métricos. Demostrar las siguientes proposiciones:
    \begin{enumerate}
        \item[(a)] Si la función \( f : X \to Y \) es Lipschitz entonces es uniformemente continua.
        \item[(b)] La función \( f : \mathbb{R} \to \mathbb{R} \) dada por \( f(t) = t^{2} \) es continua pero no es uniformemente continua.
        \item[(c)] Lo mismo para la función \( g : (0, 1) \to \mathbb{R} \) dada por \( g(t) = \frac{1}{t} \).
    \end{enumerate}
\end{ejercicio}
\begin{proof}
    Si \( f \) es Lipschitz, existe una constante \( L \ge 0 \) tal que \( d_Y(f(x_1), f(x_2)) \leq L \cdot d_X(x_1, x_2) \) para todo par de puntos en \( X \). Si \( L = 0 \), la función es constante y trivialmente uniformemente continua.

    Asumamos \( L > 0 \). Dado \( \varepsilon > 0 \), basta tomar \( \delta = \frac{\varepsilon}{L} \). Entonces, para todo \( x_1, x_2 \in X \) que cumplan \( d_X(x_1, x_2) < \delta \), vale que:
    \[
        d_Y(f(x_1), f(x_2)) \leq L \cdot d_X(x_1, x_2) < L \cdot \frac{\varepsilon}{L} = \varepsilon
    \]
    Por lo tanto, \( f \) es uniformemente continua.

    \vspace{0.3cm}
    Que la función \( f(t) = t^2 \) es continua en \( \mathbb{R} \) es un resultado conocido (producto de funciones continuas).

    Supongamos por absurdo que \( f \) es uniformemente continua. Luego, para todo \( \varepsilon > 0 \) existe \( \delta > 0 \) tal que:
    \[
        | x - y | < \delta \implies |x^2 - y^2| < \varepsilon
    \]
    En particular, esto debe valer para \( \varepsilon = 1 \). Consideremos para un \( x \in \mathbb{R} \) cualquiera el punto \( y = x + \frac{\delta}{2} \). Verificamos que la distancia en el dominio cumple:
    \[
        |x - y| = \left| x - \left( x + \frac{\delta}{2} \right) \right| = \left| -\frac{\delta}{2} \right| = \frac{\delta}{2} < \delta
    \]
    Por nuestra suposición, debería cumplirse que \( |x^2 - y^2| < 1 \). Evaluamos:
    \[
        |x^2 - y^2| = \left| x^2 - {\left(x+\frac{\delta}{2}\right)}^2 \right| = \left| x^2 - \left( x^2 + x\delta + \frac{\delta^2}{4} \right) \right| = \left| x\delta + \frac{\delta^2}{4} \right|
    \]
    Necesitamos que \( \left| x\delta + \frac{\delta^2}{4} \right| < 1 \). Pero como el dominio es \( \mathbb{R} \) y \( \delta > 0 \) está fijo, podemos elegir un \( x \) lo suficientemente grande (por ejemplo \( x > \frac{1}{\delta} \)) para que esta expresión sea estrictamente mayor a 1, violando la implicación. Por lo tanto, \( f \) no es uniformemente continua.

    \vspace{0.3cm}
    La función \( g(t) = \frac{1}{t} \) es continua en \( (0, 1) \) por ser cociente de continuas con denominador no nulo.

    Supongamos que \( g \) es uniformemente continua en \( (0,1) \). Luego, para \( \varepsilon = 1 \) debe existir un \( \delta > 0 \) tal que:
    \[
        |x - y| < \delta \implies \left| \frac{1}{x} - \frac{1}{y} \right| < 1
    \]
    Consideremos las sucesiones \( x_n = \frac{1}{n+1} \) e \( y_n = \frac{1}{n+2} \). Para todo \( n \ge 1 \), ambas sucesiones están en el intervalo \( (0, 1) \). Calculamos la distancia entre ellas:
    \[
        |x_n - y_n| = \left| \frac{1}{n+1} - \frac{1}{n+2} \right| = \frac{1}{(n+1)(n+2)}
    \]
    Como \( \lim_{n \to \infty} \frac{1}{(n+1)(n+2)} = 0 \), por definición de límite existe un \( n_0 \in \mathbb{N} \) lo suficientemente grande tal que:
    \[
        |x_{n_0} - y_{n_0}| < \delta
    \]
    Sin embargo, si evaluamos la distancia de sus imágenes, obtenemos:
    \[
        |g(x_{n_0}) - g(y_{n_0})| = | (n_0+1) - (n_0+2) | = |-1| = 1
    \]
    Para que se cumpliera la continuidad uniforme, debíamos tener que \( |g(x_{n_0}) - g(y_{n_0})| < \varepsilon = 1 \), pero hemos obtenido \( 1 < 1 \), lo cual es un absurdo. En consecuencia, \( g \) no es uniformemente continua.
\end{proof}

\begin{ejercicio}{10}
    Decidir la compacidad del conjunto \( [0, 1] \) respecto a la topología
    \[ \tau_c = \{ A : [0,1] \setminus A \text{ es a lo sumo numerable o es todo } [0, 1] \} \]
\end{ejercicio}
\begin{proof}
    El conjunto \( [0, 1] \) dotado de la topología co-numerable \( \tau_c \) \textbf{no es compacto}.

    Para demostrarlo, basta con exhibir un cubrimiento abierto que no posea ningún subcubrimiento finito. Tomemos un subconjunto infinito numerable cualquiera de \( [0, 1] \), por ejemplo, una enumeración de los racionales en dicho intervalo, o simplemente una sucesión de puntos distintos:\[
        C = \{ x_n \}_{n \in \mathbb{N}} \subseteq [0, 1]
    \]

    Para cada \( n \in \mathbb{N} \), definimos el conjunto:\[
        U_n = [0, 1] \setminus (C \setminus \{ x_n \})
    \]

    Veamos primero que cada \( U_n \) es un abierto en \( \tau_c \). En efecto, su complemento es:\[
        [0, 1] \setminus U_n = C \setminus \{ x_n \}
    \]
    Como \( C \) es numerable, cualquier subconjunto suyo también lo es. Por lo tanto, \( [0, 1] \setminus U_n \) es numerable, lo que implica que \( U_n \in \tau_c \).

    A continuación, probemos que la familia \( {\{ U_n \}}_{n \in \mathbb{N}} \) es un cubrimiento de \( [0, 1] \). Sea \( x \in [0, 1] \), tenemos dos casos:
    \begin{itemize}
        \item Si \( x \notin C \), entonces \( x \in [0, 1] \setminus C \subseteq U_1 \), por lo que está cubierto (de hecho, pertenece a todos los \( U_n \)).
        \item Si \( x \in C \), entonces existe un índice \( k \in \mathbb{N} \) tal que \( x = x_k \). Por definición, \( x_k \notin C \setminus \{ x_k \} \), lo que implica de forma directa que \( x_k \in U_k \).
    \end{itemize}
    De este modo, todo elemento de \( [0, 1] \) está en algún \( U_n \), por lo que \( [0, 1] \subseteq \bigcup_{n=1}^\infty U_n \).

    Finalmente, veamos que de este cubrimiento no se puede extraer un subcubrimiento finito. Supongamos por el absurdo que existen finitos índices \( n_1, n_2, \dots, n_k \) tales que:\[
        [0, 1] \subseteq \bigcup_{i=1}^k U_{n_i}
    \]

    Dado que tenemos infinitos elementos en la sucesión \( C \) y solo tomamos una cantidad finita de índices para el subcubrimiento, podemos elegir un índice \( m \in \mathbb{N} \) tal que \( m \neq n_i \) para todo \( i = 1, \dots, k \).

    Consideremos el punto \( x_m \in C \). Para que sea un subcubrimiento válido, \( x_m \) debería pertenecer a algún \( U_{n_i} \). Sin embargo, como \( m \neq n_i \), resulta que el punto \( x_m \) sí pertenece al conjunto \( C \setminus \{ x_{n_i} \} \). Pero esto significa que:\[
        x_m \notin [0, 1] \setminus (C \setminus \{ x_{n_i} \}) = U_{n_i}
    \]

    Como esto ocurre para todo \( i = 1, \dots, k \), concluimos que \( x_m \) no está en la unión de esos finitos abiertos. Esto contradice que sea un subcubrimiento de \( [0, 1] \).

    Por lo tanto, hemos construido un cubrimiento abierto que no admite ningún subcubrimiento finito, lo que demuestra que el espacio no es compacto.
\end{proof}

\begin{ejercicio}{11}
    \begin{enumerate}
        \item[(a)] Sea \( C \) una componente conexa de un espacio topológico Hausdorff compacto \( X \). Entonces \( C \) es la intersección de los subconjuntos clopen de \( X \) que lo contienen.
        \item[(b)] Sea \( X \) un espacio topológico conexo, compacto y de Hausdorff; \( x \in X \) y \( C \) una componente conexa de \( X - \{ x \} \), entonces \( x \in \cl{C} \).
    \end{enumerate}
\end{ejercicio}
\begin{proof}

\end{proof}

\section{Compactificaciones}

\begin{ejercicio}{12}
    Sea \( \tilde{X} \) la compactificación de Alexandrov del espacio topológico X.
    \begin{enumerate}
        \item[(a)] Tomando \( X=\R \) con la topología usual, probar que \( \tilde{X}\cong S^{1} \).
        \item[(b)] Tomando \( X=\N \) con la topología discreta, mostrar que los únicos compactos de \( \N \) son los finitos. Mostrar que \( \tilde{X} \cong \left \{ \frac{1}{n} : n \in \N \right \} \cup \{ 0 \} \) (con la topología usual de \( \R \)).
        \item[(c)] Tomando \( X=(0,1)\cup(2,3) \) con la topología usual, probar que \( \tilde{X} \) es homeomorfo a ocho en \( \R^{2} \).
        \item[(d)] Mostrar que para el caso \( X=\Q \) sucede que \( \tilde{X} \) no es ni Hausdorff, a diferencia de todos los ejemplos anteriores donde \( \tilde{X} \) resultó metrizable.
    \end{enumerate}
\end{ejercicio}
\begin{proof}
    Para el inciso (a), recordemos primero que \( \R \cong (a, b) \) para cualesquiera \( a < b \) en \( \R \) vía la función homeomorfismo dada por:
    \[ f : \R \longrightarrow (a, b), \quad f(x) = a + \frac{b - a}{\pi} \left( \arctan(x) + \frac{\pi}{2} \right) \]
    Luego, tenemos que \( S^1 \setminus \{ (1, 0) \} \cong (0, 1) \) mediante la función de homeomorfismo:
    \[ g : (0, 1) \longrightarrow S^1 \setminus \{ (1, 0) \}, \quad g(t) = e^{2\pi i t} \]
    Componiendo ambos homeomorfismos, obtenemos que \( \R \cong S^1 \setminus \{ (1, 0) \} \). Si construimos la compactificación de Alexandrov de \( \R \),
    obtenemos el espacio \( \tilde{X} = \R \cup \{ \infty \} \) con la topología \( \tau_\infty = \tau \cup \tau' \), donde \( \tau \) es la topología usual
    de \( \R \) y \[
        \tau' = \{ U \cup \infty : \R \setminus U \text{ es compacto y } U \in \tau \}
    \]
    Queremos mostrar que \( \tilde{X} \cong S^1 \). Para ello, definamos la función:
    \[ h : \tilde{X} \longrightarrow S^1, \quad h(x) = \begin{cases}
            g(f(x)) & \text{si } x \in \R   \\
            (1, 0)  & \text{si } x = \infty
        \end{cases} \]
    Es fácil verificar que \( h \) es un homeomorfismo. En efecto, \( h \) es claramente biyectiva, ya que \( f \) y \( g \) lo son.
    Para mostrar que \( h \) es continua, basta con verificar la continuidad en el punto \( \infty \), ya que en el resto de los puntos es la composición
    de funciones continuas. Dado un entorno abierto \( V \) de \( (1, 0) \) en \( S^1 \), su preimagen bajo \( h \) es:
    \[ h^{-1}(V) = \{ x \in \R : g(f(x)) \in V \} \cup \{ \infty \} \]
    Dado que \( g \) es un homeomorfismo, \( g^{-1}(V) \) es un entorno abierto de \( 0 \) en \( (0, 1) \). Luego
    \[ f^{-1}(g^{-1}(V)) \] es un entorno abierto de \( \infty \) en \( \R \). Además, el complemento de \( f^{-1}(g^{-1}(V)) \) es compacto, ya que
    es homeomorfo a un subconjunto compacto de \( (0, 1) \). Por lo tanto, \( h^{-1}(V) \) es un entorno abierto de \( \infty \) en \( \tilde{X} \), lo que demuestra la continuidad de \( h \).
    Para mostrar que \( h^{-1} \) es continua, basta con verificar la continuidad en el punto \( (1, 0) \). Dado un entorno abierto \( U \) de \( \infty \) en \( \tilde{X} \), su preimagen bajo \( h^{-1} \) es:
    \[ h^{-1}(U) = \{ z \in S^1 : h^{-1}(z) \in U \} \]
    Dado que \( h^{-1} \) es un homeomorfismo, \( h^{-1}(U) \) es un entorno abierto de \( (1, 0) \) en \( S^1 \). Esto demuestra la continuidad de \( h^{-1} \). Por lo tanto, \( h \) es un homeomorfismo entre
    \( \tilde{X} \) y \( S^1 \).

    Para el iniciso (b), si \( A \subseteq \N \) es un conjunto finito, para cualquier cubrimiento de abiertos \( \rho \) de \( A \) basta con tomar
    para cada \( a \in A \) un abierto \( U_a \in \rho \) tal que \( a \in U_a \). Entonces, la familia \( \{ U_a : a \in A \} \) es un
    subcubrimiento finito de \( A \implies A \) es compacto. Si \( A \) es infinito, consideremos la familia de abiertos \( \{ \{ n \} : n \in A \} \).
    Esta familia es un cubrimiento de \( A \) que no admite ningún subcubrimiento finito, ya que cada abierto solo cubre un punto y \( A \) tiene infinitos.
    Por lo tanto, los únicos compactos de \( \N \) son los finitos. Veamos ahora que \( \tilde{X} \cong \{ \frac{1}{n} : n \in \N \} \cup \{ 0 \} \).
    Para ello, definamos la función \( f : \N \to \{ \frac{1}{n} : n \in \N \} \) dada por:\[
        f(n) = \frac{1}{n}
    \]
    Sea \( \tau \) la topología discreta y \( \tau' = \{ U \cup \infty : U^c \text{ es finito y } U \in \tau \} \) la topología de los entornos de \( \infty \)
    entonces \( \tau_\infty = \tau \cup \tau' \) es la topología de \( \tilde{X} \). Definamos la extensión de \( f \) a \( \tilde{X} \) por:
    \[
        \tilde{f} : \tilde{X} \longrightarrow \left \{ \frac{1}{n} : n \in \N \right \} \cup \{ 0 \}, \quad \tilde{f}(x) = \begin{cases}
            f(x) & \text{si } x \in \N   \\
            0    & \text{si } x = \infty
        \end{cases}
    \]
    Análogamente, \( \tilde{f} \) es biyectiva y para mostrar que es un homeomorfismo basta con verificar la continuidad en el punto \( \infty \).
    Dado un entorno abierto \( V \) de \( 0 \) en \( \{ \frac{1}{n} : n \in \N \} \cup \{ 0 \} \), existe un \( N \in \N \) tal que \( V \) contiene a todos los puntos \( \frac{1}{n} \) con \( n > N \).
    Por lo tanto, el complemento de \( \tilde{f}^{-1}(V) \) es finito, ya que solo puede contener a los puntos \( 1, \frac{1}{2}, \dots, \frac{1}{N} \).
    Esto implica que \( \tilde{f}^{-1}(V) \) es un entorno abierto de \( \infty \) en \( \tilde{X} \), lo que demuestra la continuidad de \( \tilde{f} \).
    De manera similar, la continuidad de \( \tilde{f}^{-1} \) se verifica en el punto \( 0 \). Por lo tanto, \( \tilde{f} \) es un homeomorfismo entre \( \tilde{X} \) y \( \{ \frac{1}{n} : n \in \N \} \cup \{ 0 \} \).

    Para el inciso (c), notemos que \( X = (0, 1) \cup (2, 3) \) está formado por dos componentes conexas.
    Podemos ver al ocho geométrico como el subespacio de \(\R^2\) formado por dos circunferencias tangentes en el origen \((0,0)\).
    Llamemos \(C_1\) a la circunferencia superior y \(C_2\) a la inferior. Sabemos que: \( (0,1) \cong C_1 \setminus \{(0,0)\} \) y \((2,3) \cong C_2 \setminus \{(0,0)\} \).
    Por lo tanto, existe un homeomorfismo \(f: X \to (C_1 \cup C_2) \setminus \{(0,0)\} \).
    Al construir la compactificación de Alexandrov \(\tilde{X}\), le agregamos el punto \(\infty \). Definimos la función \(h: \tilde{X} \to C_1 \cup C_2\) (el ocho) como \(h(x) = f(x)\) si \(x \in X\), y \(h(\infty) = (0,0)\).
    Dado que el complemento de cualquier entorno abierto del \((0,0)\) en el ocho es un conjunto compacto que no contiene al origen (son porciones cerradas de las circunferencias), su preimagen bajo \(f^{-1}\) será un compacto en \(X\).
    Esto garantiza que los entornos de \(\infty \) en \(\tilde{X}\) se corresponden con los entornos de \((0,0)\) en el ocho, haciendo de \(h\) un homeomorfismo.

    Para el inciso (d), supongamos por el absurdo que \( \tilde{X} \) es Hausdorff. Entonces, dado un punto \( p \in \Q \),
    existen entornos abiertos disjuntos \( U \) y \( V \) en \( \tilde{X} \) tales que \( p \in U \) y \( \infty \in V \).
    Como \( V \) es un entorno abierto de \( \infty \), por definición de la topología de Alexandrov, \( \Q \setminus V \)
    debe ser un conjunto compacto en \( \Q \).
    Como \( U \) y \( V \) son disjuntos, tenemos que \( U \subseteq \Q \setminus V \). Al ser \( U \) un abierto en \( \tilde{X} \) que no contiene a
    \( \infty \), \( U \) es un abierto no vacío en \( \Q \).
    Esto implicaría que el compacto \(\Q \setminus V\) contiene un conjunto abierto no vacío de \(\Q \). Sin embargo, \( {(\Q \setminus V)}^\circ = \varnothing \implies U^\circ = U \subseteq \varnothing \) absurdo.
    Por lo tanto, ningún compacto de \(\Q \) puede contener un abierto no vacío de \( \Q \).
    Llegamos a una contradicción. En conclusión, \( \tilde{X} \) no es Hausdorff.
\end{proof}

\begin{ejercicio}{13}
    Mostrar que si dos espacios topológicos son homeomorfos entonces también son homeomorfos sus compactificaciones de Alexandrov.
\end{ejercicio}
\begin{proof}
    Sea \( X \) un espacio topológico y \( \tilde{X} = X \cup \{ \infty \} \) su compactificación de Alexandrov. Supongamos que \( Y \) es otro espacio topológico tal que \( X \cong Y \) mediante un homeomorfismo \( f : X \to Y \).

    Queremos mostrar que \( \tilde{X} \cong \tilde{Y} \). Para ello, definamos la extensión \( \tilde{f} : \tilde{X} \to \tilde{Y} \) por:
    \[
        \tilde{f}(x) = \begin{cases}
            f(x)    & \text{si } x \in X    \\
            \infty' & \text{si } x = \infty
        \end{cases}
    \]
    Donde \( \infty' \) es el punto agregado en \( \tilde{Y} \). Es evidente que \( \tilde{f} \) es biyectiva ya que \( f \) es biyectiva y \( \tilde{f}(\infty) = \infty' \).

    Probemos que \( \tilde{f} \) es continua. Sea \( V \) un abierto en \( \tilde{Y} \). Por la definición de la topología de Alexandrov, tenemos dos casos:
    \begin{itemize}
        \item \textbf{Caso 1:} \( \infty' \notin V \). Entonces \( V \) es un abierto de \( Y \). Su preimagen es \( \tilde{f}^{-1}(V) = f^{-1}(V) \). Como \( f \) es continua, \( f^{-1}(V) \) es abierto en \( X \), y por ende también es abierto en \( \tilde{X} \).
        \item \textbf{Caso 2:} \( \infty' \in V \). Entonces \( V = W \cup \{ \infty' \} \), donde \( W \) es un abierto de \( Y \) tal que \( Y \setminus W \) es un compacto en \( Y \). Su preimagen bajo \( \tilde{f} \) es:
              \[
                  \tilde{f}^{-1}(V) = f^{-1}(W) \cup \{ \infty \}
              \]
              Como \( f \) es un homeomorfismo, \( f^{-1}(W) \) es un abierto de \( X \). Además, su complemento en \( X \) cumple que:
              \[
                  X \setminus f^{-1}(W) = f^{-1}(Y \setminus W)
              \]
              Como \( Y \setminus W \) es compacto en \( Y \) y \( f^{-1} \) es continua, la imagen inversa \( f^{-1}(Y \setminus W) \) es compacta en \( X \). Por lo tanto, por definición de la topología de Alexandrov, \( \tilde{f}^{-1}(V) \) es un abierto en \( \tilde{X} \).
    \end{itemize}
    Esto demuestra que \( \tilde{f} \) es continua.

    Para concluir que es un homeomorfismo, veamos de manera análoga que \( \tilde{f} \) es una función abierta. Sea \( U \) un abierto en \( \tilde{X} \):
    \begin{itemize}
        \item \textbf{Caso 1:} \( \infty \notin U \). Entonces \( U \) es un abierto en \( X \). Su imagen es \( \tilde{f}(U) = f(U) \), que es abierta en \( Y \) (y por tanto en \( \tilde{Y} \)) porque \( f \) es una función abierta.
        \item \textbf{Caso 2:} \( \infty \in U \). Entonces \( U = A \cup \{ \infty \} \), donde \( A \) es un abierto de \( X \) tal que \( X \setminus A \) es compacto en \( X \). Su imagen es:
              \[
                  \tilde{f}(U) = f(A) \cup \{ \infty' \}
              \]
              Como \( f \) es homeomorfismo, \( f(A) \) es abierto en \( Y \). Evaluando el complemento en \( Y \) obtenemos:
              \[
                  Y \setminus f(A) = f(X \setminus A)
              \]
              Puesto que \( X \setminus A \) es compacto en \( X \) y \( f \) es continua, \( f(X \setminus A) \) es un compacto en \( Y \). Esto significa que \( \tilde{f}(U) \) cumple con las condiciones para ser un abierto de \( \tilde{Y} \).
    \end{itemize}
    Al ser \( \tilde{f} \) biyectiva, continua y abierta, concluimos que \( \tilde{f} : \tilde{X} \to \tilde{Y} \) es un homeomorfismo, por lo que \( \tilde{X} \cong \tilde{Y} \).
\end{proof}

\begin{ejercicio}{14}
    Sea \( X \) un espacio topológico de Tychonoff y sea \( RU(X) \) el conjunto de todas sus redes universales.
    \begin{enumerate}
        \item[(a)] Dada \( x=(x_{i})_{i\in I}\in RU(X) \) e \( f\in C_{b}(X) \), probar que existe un \( x_{f}\in\mathbb{C} \) tal que \( f(x_{i})\longrightarrow_{i\in I}x_{f} \).
        \item[(b)] Sean \( \mathcal{F}_{1} \), \( B(X) \) y \( F \) según la construcción que se hace en los apuntes. Mostrar que \( F \) separa puntos de \( B(X) \).
        \item[(c)] Probar que \( (B(X),\tau_{\mathcal{F}}) \) es Hausdorff, donde \( \tau_{\mathcal{F}} \) es la topología inicial dada por \( F \).
        \item[(d)] Probar que la función \( G:X\rightarrow B(X) \) dada por \( G(x)=\overline{c_{x}} \), donde \( c_{x} \) es la sucesión constante igual a \( x \), es un embedding.
        \item[(e)] Probar que \( (B(X),\tau_{\mathcal{F}}) \) es un compacto Hausdorff y que el par \( (B(X),G) \) es una H-compactificación de \( X \) que es isomorfa a la de Stone-Čech.
    \end{enumerate}
\end{ejercicio}
\begin{proof}

\end{proof}

\begin{ejercicio}{15}
    Sea \( X \) discreto y \( (F,\beta(X)) \) su compactificación de Stone-Čech.
    \begin{enumerate}
        \item[(a)] Para \( A\subseteq X \), entonces \( \overline{F(A)} \) y \( \overline{F(A^{C})} \) son disjuntos.
        \item[(b)] Si \( U \) es abierto en \( \beta(X) \), entonces \( \overline{U} \) es abierto de \( \beta(X) \).
        \item[(c)] \( \beta(X) \) es totalmente disconexo.
    \end{enumerate}
\end{ejercicio}
\begin{proof}

\end{proof}

\begin{ejercicio}{16}
    Probar que un espacio Tychonoff es conexo \( \iff \) su compactificación de Stone-Čech es conexa.
\end{ejercicio}
\begin{proof}

\end{proof}

\end{document}